\documentclass{exam}
\usepackage{../../commonheader}

\discnumber{8}
\title{Mutation, Classes, \& Attributes}
\date{October 12, 2026 -- October 16, 2026}

% Exercises and lessons are imported from the topic bank when possible.
\begin{document}
\maketitle

\begin{meta}
\small
\textbf{Example Timeline}
\begin{itemize}
    \item \textbf{This is a menu of problems, not a requirement to finish the worksheet.} Choose or skip problems based on your students, and give each mini-lecture immediately before its section.
    \item Check-in and icebreaker [5 min]. Ask students for one operation that constructs a list and one that mutates a list.
    \item List mutation: selected parts of Q1 (WWPD), then Q2 (environment diagram) [15 min]. Track shallow copies, shared nested lists, and the return values of mutating methods. Q3, \lstinline{insert_n}, is extra implementation practice.
    \item Classes and dataclasses: compare the two \lstinline{Car} definitions, then selected parts of Q4 (museum tickets) [15 min]. Connect constructor arguments and defaults to instance attributes, then trace method calls, aliasing, and reassignment.
    \item Attributes and methods: selected parts of Q5 (\lstinline{Account}) [10 min]. Distinguish functions from bound methods and trace dot expressions, attribute lookup, and assignment.
    \item Choose Q6 (counters) or Q7 (labels at a depth) [10 min]. Q6 practices saved bound methods and reassignment; draw the class and both instances. Q7 practices tree recursion and mutating a result list while leaving the input tree unchanged.
    \item Wrap-up and questions [5 min]. Ask students to explain mutation, attribute lookup, and bound methods in their own words.
\end{itemize}
\textbf{Teaching Goals}
\begin{itemize}
    \item Distinguish construction, reassignment, aliasing, and mutation. Mutation changes an existing object, so every reference to that object observes the change.
    \item Use box-and-pointer or environment diagrams to track list mutation, shallow copies, and shared nested objects. For instance mutation, draw attributes inside the instance that owns them.
    \item In Q7, each recursive call returns a list of labels at the requested depth. Use \lstinline{extend} to combine those results, and return the list separately because \lstinline{extend} returns \lstinline{None}. Only the result lists are mutated, not the tree.
    \item Recognize that classes and instances are objects with attributes. Functions and bound methods are objects as well.
    \item Compare explicit and generated initializers for \lstinline{Car}: both accept \lstinline{owner} and optional \lstinline{gas}; Python supplies \lstinline{self}. The unannotated \lstinline{wheels} remains a class attribute.
    \item Dataclasses support methods with parameters, return values, and mutation. In Q4, distinguish the boolean returned by \lstinline{admit} from its change to \lstinline{remaining}. Equal field values do not imply object identity.
    \item For Q5's ordinary attributes, lookup checks the instance and then its class. Assignment writes to the object before the dot: \lstinline{a.interest = 7} shadows the class attribute; \lstinline{Account.interest = 4} changes it.
    \item A function stored as a class attribute acts as a method. Looking it up through an instance produces a bound method: \lstinline{a.deposit(5)} is equivalent to \lstinline{Account.deposit(a, 5)}. In Q6, a saved bound method retains its instance even after the original name is reassigned.
\end{itemize}
\end{meta}

\section{List Mutation and Environment Diagrams}
\begingroup
\subimport{../../topics/lists/mutable/text/}{mutation_overview.tex}
\endgroup
\begin{questions}
    \subimport{../../topics/lists/mutable/easy/}{wwpd-mixed.tex}

    \subimport{../../topics/lists/mutable/easy/env-diagram/}{mystery.tex}

    \subimport{../../topics/lists/mutable/medium/}{insert-n.tex}
\end{questions}

\newpage
\section{Classes and Instances}
\begingroup
\subimport{../../topics/oop/text/}{oop_overview.tex}
\endgroup

\subsection*{The Same Data with a Dataclass}
\begingroup
\subimport{../../topics/oop/dataclasses/text/}{dataclass.tex}
\endgroup

\begin{questions}
    \setcounter{question}{3}
    \subimport{../../topics/oop/dataclasses/easy/}{ticket-mutation.tex}
\end{questions}

\newpage
\section{Attribute Lookup, Assignment, and Bound Methods}
\begin{questions}
\setcounter{question}{4}
    \subimport{../../topics/oop/easy/wwpd/}{account.tex}

    \subimport{../../topics/oop/easy/wwpd/}{counter-attributes.tex}
\end{questions}

\newpage
\section{Tree Recursion and Result Lists}
\begin{blocksection}
\subimport{../../topics/trees/class/text/}{implementation.tex}
\end{blocksection}
\begin{questions}
    \setcounter{question}{6}
    \subimport{../../topics/trees/class/medium/}{labels-at-depth.tex}
\end{questions}

\end{document}
